\documentclass[11pt]{article}
\usepackage{explorer}
\newcommand{\ExplorerLanguage}{EN}
\newcommand{\ExplorerNumber}{12}
\newcommand{\ExplorerKind}{Worked solutions}
\begin{document}
\ExplorerTitle{Upper high school · analysis and algorithms}{Bracketing a root by bisection}
\ExplorerSection{Exercise data}
For exercise 1: seek a zero of \(f(x)=x^{2}-2\) on [1,2]. After two halvings, the retained interval is [1.25,1.5].\par
\ExplorerSection{Exercise 1}
Carry out a third step of the example.\par
The midpoint is 1.375, whose square is \(1.890625<2\). The new interval is [1.375,1.5].\par
\ExplorerSection{Exercise 2}
Starting with width 1, how many halvings guarantee width at most \(\frac{1}{1000}\)?\par
\(2^{9}=512<1000\leq 1024=2^{10}\), so 10 halvings suffice and 9 do not.\par
\ExplorerSection{Exercise 3}
Why does the sign change of \(f(x)=\frac{1}{x}\) between −1 and 1 not prove the existence of a zero?\par
The function is undefined at 0 and is not continuous on [−1,1]. The theorem’s hypothesis fails; \(\frac{1}{x}\) is never zero.\par
\ExplorerSection{Exercise 4}
A root lies in [1.4,1.5]. What maximum error is guaranteed by midpoint 1.45?\par
The half-width is \(\frac{1.5-1.4}{2}=0.05\).\par
\ExplorerSection{Problem: complete reasoning}
For \(f(x)=x^{3}-x-1\) on [1,2], prove existence and uniqueness of a root \(\alpha \). Carry out three bisection steps. How many halvings guarantee that the final midpoint approximates \(\alpha \) with error at most 0.001?\par
\(f(1)=-1\) and \(f(2)=5\), and f is continuous. Also \(f'(x)=3x^{2}-1\geq 2\) on [1,2], so f is strictly increasing. Midpoints 1.5, 1.25, 1.375 give [1,1.5], [1.25,1.5], [1.25,1.375] because f takes values \(0.875, -0.296875\) and 0.224609375. After n halvings, midpoint error is at most \(\frac{1}{2^{n+1}}\). Thus \(2^{n+1}\geq 1000\), requiring \(n\geq 9\).\par
\end{document}
